\chapter{Testing}

\section{Testing Strategy}
Testing was performed at three levels: unit testing of engine modules, integration testing of the full simulation pipeline, and performance testing under target hardware constraints.

\section{Unit Testing -- Terrain Engine}

\begin{table}[h]
\centering
\caption{Terrain Engine Test Cases}
\label{tab:test_terrain}
\begin{tabular}{|c|p{4cm}|p{4cm}|p{3cm}|}
\hline
\textbf{TC-ID} & \textbf{Test Case} & \textbf{Expected Result} & \textbf{Status} \\
\hline
TC-01 & Generate 256$\times$256 heightmap & 65,536 float values, no NaN & \textcolor{green!60!black}{PASS} \\
\hline
TC-02 & Same seed produces same terrain & Identical heightmap arrays & \textcolor{green!60!black}{PASS} \\
\hline
TC-03 & Different seeds produce different terrain & Heightmaps differ & \textcolor{green!60!black}{PASS} \\
\hline
TC-04 & Height range is bounded & $h \in [-30, +30]$ approx. & \textcolor{green!60!black}{PASS} \\
\hline
TC-05 & Crater at known position & Local minimum at crater center & \textcolor{green!60!black}{PASS} \\
\hline
TC-06 & Bilinear sampling accuracy & Interpolated height matches mesh & \textcolor{green!60!black}{PASS} \\
\hline
TC-07 & Slope computation at flat region & Slope $\approx 0$ & \textcolor{green!60!black}{PASS} \\
\hline
TC-08 & Slope computation at crater rim & Slope $> 0.4$ & \textcolor{green!60!black}{PASS} \\
\hline
\end{tabular}
\end{table}

\section{Unit Testing -- Physics Engine}

\begin{table}[h]
\centering
\caption{Physics Engine Test Cases}
\label{tab:test_physics}
\begin{tabular}{|c|p{4cm}|p{4cm}|p{3cm}|}
\hline
\textbf{TC-ID} & \textbf{Test Case} & \textbf{Expected Result} & \textbf{Status} \\
\hline
TC-09 & Free fall from 100m (no thrust) & Impact at $v = \sqrt{2 \cdot 3.72 \cdot 100} \approx 27.3$ m/s & \textcolor{green!60!black}{PASS} \\
\hline
TC-10 & Hover at constant altitude & Throttle $= g/F_{\max} = 0.31$ & \textcolor{green!60!black}{PASS} \\
\hline
TC-11 & Fuel depletion at full throttle & Fuel reaches 0 at $t = 125$s & \textcolor{green!60!black}{PASS} \\
\hline
TC-12 & Collision detection at ground & Lander stops at groundHeight + 1.5 & \textcolor{green!60!black}{PASS} \\
\hline
TC-13 & Soft landing classification & $|v_y| < 3$ m/s $\Rightarrow$ landed=true & \textcolor{green!60!black}{PASS} \\
\hline
TC-14 & Crash classification & $|v_y| > 3$ m/s $\Rightarrow$ crashed=true & \textcolor{green!60!black}{PASS} \\
\hline
TC-15 & Lateral drag reduces velocity & $v_x$ decays toward zero over time & \textcolor{green!60!black}{PASS} \\
\hline
\end{tabular}
\end{table}

\section{Integration Testing -- Autopilot}

The autopilot was tested across 50 simulated landings with varying initial positions over the terrain. Results:

\begin{table}[h]
\centering
\caption{Autopilot Performance Summary (50 Landings)}
\label{tab:autopilot_results}
\begin{tabular}{|p{4cm}|c|}
\hline
\textbf{Metric} & \textbf{Value} \\
\hline
Successful landings & 48 / 50 (96\%) \\
\hline
Mean impact speed & 1.2 m/s \\
\hline
Max impact speed (successful) & 2.7 m/s \\
\hline
Mean fuel remaining & 89.3 kg / 100 kg \\
\hline
Mean mission time & 31.4 s \\
\hline
Crash causes & Steep crater rim (2 cases) \\
\hline
\end{tabular}
\end{table}

\section{Performance Testing}

\begin{table}[h]
\centering
\caption{Performance Benchmarks}
\label{tab:perf}
\begin{tabular}{|p{4cm}|c|c|c|}
\hline
\textbf{Metric} & \textbf{Target} & \textbf{Measured} & \textbf{Status} \\
\hline
Frame rate (orbital) & $\geq$ 30 FPS & 55--60 FPS & \textcolor{green!60!black}{PASS} \\
\hline
Frame rate (descent) & $\geq$ 30 FPS & 50--58 FPS & \textcolor{green!60!black}{PASS} \\
\hline
Initial load time & $< 3$s & 1.8s & \textcolor{green!60!black}{PASS} \\
\hline
Terrain generation & $< 500$ms & 120ms & \textcolor{green!60!black}{PASS} \\
\hline
GPU memory usage & $< 2$~GB & $\sim$800 MB & \textcolor{green!60!black}{PASS} \\
\hline
JS heap size & $< 200$ MB & $\sim$85 MB & \textcolor{green!60!black}{PASS} \\
\hline
Bundle size (prod) & $< 5$ MB & 2.1 MB & \textcolor{green!60!black}{PASS} \\
\hline
\end{tabular}
\end{table}

\textit{Hardware:} Windows laptop, NVIDIA GTX 1650 (4~GB VRAM), 16~GB RAM, Chrome 120.

\section{Cross-Browser Testing}

\begin{table}[h]
\centering
\caption{Browser Compatibility}
\begin{tabular}{|p{3cm}|c|c|c|}
\hline
\textbf{Feature} & \textbf{Chrome 120} & \textbf{Firefox 121} & \textbf{Edge 120} \\
\hline
WebGL 2.0 & \checkmark & \checkmark & \checkmark \\
\hline
Post-processing & \checkmark & \checkmark & \checkmark \\
\hline
Web Audio API & \checkmark & \checkmark & \checkmark \\
\hline
OrbitControls & \checkmark & \checkmark & \checkmark \\
\hline
CSS backdrop-filter & \checkmark & \checkmark & \checkmark \\
\hline
\end{tabular}
\end{table}

\section{Discussion}
All functional requirements (FR-01 through FR-11) and non-functional requirements (NFR-01 through NFR-06) were met. The autopilot achieves a 96\% success rate, with the two failures occurring when the selected landing zone was on the steep inner wall of a crater---a scenario that a real mission would also avoid during orbital survey. The system maintains well above the target 30 FPS on mid-range hardware.
